\noindent
Recall that given a scheme $X$, an invertible sheaf $\mathcal{L}$
on $X$, and a sheaf of $\mathcal{O}_X$-modules $\mathcal{F}$
we get a graded ring
$\Gamma_*(X, \mathcal{L}) =
\bigoplus\nolimits_{n \geq 0} \Gamma(X, \mathcal{L}^{\otimes n})$
and a graded $\Gamma_*(X, \mathcal{L})$-module
$\Gamma_*(X, \mathcal{L}, \mathcal{F}) =
\bigoplus\nolimits_{n \in \mathbf{Z}}
\Gamma(X, \mathcal{F} \otimes_{\mathcal{O}_X} \mathcal{L}^{\otimes n})$
see Modules, Definition \ref{modules-definition-gamma-star}.
If we have moreover a section $s \in \Gamma(X, \mathcal{L})$, then
we obtain a map
\begin{equation}
\label{equation-module-invert-s}
\Gamma_*(X, \mathcal{L}, \mathcal{F})_{(s)}
\longrightarrow
\Gamma(X_s, \mathcal{F}|_{X_s})
\end{equation}
which sends $t/s^n$ where
$t \in \Gamma(X, \mathcal{F} \otimes_{\mathcal{O}_X} \mathcal{L}^{\otimes n})$
to $t|_{X_s} \otimes s|_{X_s}^{-n}$. This makes sense
because $X_s \subset X$ is by definition the open over which
